%source mathadata
\begin{multicols}{2}
		À chaque lancer, on regarde de quelle cible il est le plus proche : ce lancer est alors attribué à la joueuse correspondante. Cette joueuse marque un point. En cas d'égalité, aucun point n'est attribué. Au bout de 6 lancers, la joueuse ayant le plus de points a gagné.

		\begin{center}
			\includegraphics[scale=0.8]{images/dart2.png}
		\end{center}
	\end{multicols}

	\begin{multicols}{2}

		\begin{MADenumerate}

			\item Retracer les points cibles $S$ et $A$ sur le nouveau graphique.

			\item Tracer chaque lancer dans le repère, puis compléter le tableau ci-dessous. Pour les lancers marqués d'un $*$, calculer les deux distances avant de conclure.

			\begin{center}
				\begin{tabular}{|c|c|c|c|}
					\hline
					\rowcolor{floralwhite}
					Lancer & Joueuse & Coordonnées & Point attribué à \\ \hline
					1$*$   & Samia   & $(2;3)$     &                  \\ \hline
					2      & Samia   & $(5;3)$     &                  \\ \hline
					3$*$   & Alicia  & $(1;5)$     &                  \\ \hline
					4$*$   & Samia   & $(2;4)$     &                  \\ \hline
					5      & Alicia  & $(5;4)$     &                  \\ \hline
					6      & Samia   & $(0;3)$     &                  \\ \hline
				\end{tabular}
			\end{center}

			\item Qui a gagné la partie ?

		\end{MADenumerate}

		\columnbreak

		\begin{figure}[H]
			\centering
			\begin{tikzpicture}[line cap=round,line join=round,>=triangle 45,x=1cm,y=1cm,scale=0.8]
				\begin{axis}[
						x=1cm,y=1cm,
						axis lines=middle,
						ymajorgrids=true,
						xmajorgrids=true,
						xmin=-1,
						xmax=7,
						ymin=-1,
						ymax=6,
						xtick={-1,...,7},
						ytick={-1,...,5},
					]
				\end{axis}
			\end{tikzpicture}
		\end{figure}

	\end{multicols}
